% =============================================================================
%  python_cours.tex - Les bases de Python pour resoudre des problemes
%  Des types de variables jusqu'aux decorateurs, avec la gestion des erreurs.
%
%  Compilation (deux passes pour la table des matieres) :
%     pdflatex python_cours.tex
%     pdflatex python_cours.tex
% =============================================================================
\documentclass[11pt,a4paper]{article}

\usepackage[utf8]{inputenc}
\usepackage[T1]{fontenc}
\usepackage[french]{babel}
\usepackage[margin=2.2cm]{geometry}
\usepackage[table]{xcolor}
\usepackage{listings}
\usepackage{tcolorbox}
\usepackage{enumitem}
\usepackage{titlesec}
\usepackage{longtable}
\usepackage{array}
\usepackage{fancyhdr}
\usepackage{hyperref}

% --- Couleurs ---------------------------------------------------------------
\definecolor{primary}{HTML}{20603D}   % vert Python fonce
\definecolor{accent}{HTML}{2E86DE}    % bleu Python
\definecolor{codebg}{HTML}{F4F7FB}
\definecolor{codekw}{HTML}{20603D}
\definecolor{codestr}{HTML}{B7472A}
\definecolor{codecom}{HTML}{7F8C99}
\definecolor{rulecol}{HTML}{D5DEE8}
\definecolor{warnbg}{HTML}{FFF4E6}
\definecolor{warnframe}{HTML}{E67E22}

\hypersetup{colorlinks=true, linkcolor=accent, urlcolor=accent, pdftitle={Cours Python - les bases}}

% --- En-tete / pied de page -------------------------------------------------
\pagestyle{fancy}
\fancyhf{}
\fancyhead[L]{\small\color{primary}Cours Python -- les bases}
\fancyhead[R]{\small\color{codecom}\thepage}
\renewcommand{\headrulewidth}{0.4pt}

% --- Titres -----------------------------------------------------------------
\titleformat{\section}{\Large\bfseries\color{primary}}{\thesection.}{0.6em}{}
\titleformat{\subsection}{\large\bfseries\color{accent}}{\thesubsection}{0.6em}{}

% --- Style du code Python ---------------------------------------------------
\lstdefinestyle{py}{
  language=Python,
  backgroundcolor=\color{codebg},
  basicstyle=\ttfamily\small,
  keywordstyle=\color{codekw}\bfseries,
  stringstyle=\color{codestr},
  commentstyle=\color{codecom}\itshape,
  numbers=none,
  frame=single,
  rulecolor=\color{rulecol},
  breaklines=true,
  showstringspaces=false,
  columns=fullflexible,
  keepspaces=true,
  aboveskip=8pt, belowskip=8pt,
  xleftmargin=6pt, xrightmargin=6pt,
  literate=%
    {é}{{\'e}}1 {è}{{\`e}}1 {ê}{{\^e}}1 {à}{{\`a}}1 {ô}{{\^o}}1
    {î}{{\^i}}1 {ç}{{\c c}}1 {É}{{\'E}}1 {È}{{\`E}}1 {ù}{{\`u}}1 {û}{{\^u}}1,
}
\lstset{style=py}

% --- Boites -----------------------------------------------------------------
\newtcolorbox{infobox}[1]{colback=codebg,colframe=accent,title=#1,
  fonttitle=\bfseries\color{white},boxrule=0.5pt,arc=2pt,left=6pt,right=6pt,top=4pt,bottom=4pt}

\newtcolorbox{piege}{colback=warnbg,colframe=warnframe,title=\textbf{Piege courant},
  fonttitle=\bfseries\color{white},boxrule=0.5pt,arc=2pt,left=6pt,right=6pt,top=4pt,bottom=4pt}

% =============================================================================
\begin{document}

% --- Page de titre ----------------------------------------------------------
\begin{titlepage}
\centering
\vspace*{3cm}
{\color{primary}\Huge\bfseries Les bases de Python\par}
\vspace{0.6cm}
{\color{accent}\LARGE Des variables aux decorateurs\par}
\vspace{1.2cm}
{\large Tout ce qu'il faut pour resoudre des problemes\\
et comprendre les erreurs.\par}
\vfill
{\color{codecom}\large Cours a imprimer -- support d'entrainement\par}
\vspace{2cm}
\end{titlepage}

\tableofcontents
\newpage

% =============================================================================
\section{Types de variables}
% =============================================================================
Python est \emph{dynamiquement type} : une variable n'a pas de type declare, elle
recoit le type de la valeur qu'on lui affecte. Elle est aussi \emph{fortement
typee} : \texttt{"3" + 5} leve une erreur au lieu de deviner.

\subsection{Les types de base}
\begin{lstlisting}
n   = 42            # int   : entier de taille illimitee
x   = 3.14          # float : nombre a virgule flottante
b   = True          # bool  : True / False (sous-classe de int)
s   = "bonjour"     # str   : chaine de caracteres (immuable)
z   = None          # NoneType : absence de valeur
c   = 2 + 3j        # complex

print(type(n))      # <class 'int'>
print(isinstance(b, int))   # True : bool herite de int
\end{lstlisting}

Chaque valeur en Python est un \emph{objet} qui connaît son propre type. Le type
determine les operations autorisees : on peut multiplier deux \texttt{int}, mais
pas additionner un \texttt{int} et une \texttt{str}. Retenir la difference entre
type \emph{mutable} (modifiable apres creation : liste, dict, set) et type
\emph{immuable} (fige : int, float, str, tuple) est essentiel : cela explique
beaucoup d'erreurs et de comportements surprenants que l'on verra plus loin.

\subsection{Conversions (casting)}
Convertir, c'est demander explicitement a Python de fabriquer un nouvel objet
d'un autre type a partir d'une valeur. C'est indispensable, par exemple, quand on
lit une donnee au clavier (\texttt{input} renvoie toujours une \texttt{str}) et
qu'on veut faire un calcul dessus. Attention : une conversion peut echouer et
lever une exception (\texttt{int("abc")} $\rightarrow$ \texttt{ValueError}).

\begin{lstlisting}
int("42")       # 42     (str  -> int)
float("3.14")   # 3.14
str(42)         # "42"
int(3.9)        # 3      (troncature, pas d'arrondi)
bool(0)         # False  ; bool("") -> False ; bool([]) -> False
\end{lstlisting}

\begin{piege}
La division \texttt{/} renvoie toujours un \texttt{float} (\texttt{6 / 2 == 3.0}).
Pour un entier, utilise la division entiere \texttt{//} (\texttt{7 // 2 == 3}).
\end{piege}

% =============================================================================
\section{Operateurs}
% =============================================================================
Les operateurs sont les symboles qui combinent ou comparent des valeurs. Deux
familles reviennent sans cesse dans la resolution de problemes : les operateurs
de \emph{comparaison} (qui renvoient un booleen \texttt{True}/\texttt{False} et
pilotent les conditions) et les operateurs \emph{logiques} (\texttt{and},
\texttt{or}, \texttt{not}) qui les combinent. Bon a savoir : \texttt{and} et
\texttt{or} sont \emph{paresseux} (court-circuit) -- \texttt{a and b} n'evalue
\texttt{b} que si \texttt{a} est vrai, ce qui permet d'ecrire des tests surs
comme \texttt{if x and x[0] == 1}.

\begin{longtable}{>{\raggedright\arraybackslash}p{4cm} p{9cm}}
\rowcolor{primary}
\textbf{\color{white}Categorie} & \textbf{\color{white}Operateurs} \\
\hline
Arithmetiques   & \texttt{+  -  *  /  //  \%  **} \\
Comparaison     & \texttt{==  !=  <  >  <=  >=} \\
Logiques        & \texttt{and  or  not} \\
Appartenance    & \texttt{in  ,  not in} \\
Identite        & \texttt{is  ,  is not} \\
Affectation     & \texttt{=  +=  -=  *=  /=  //=} \\
\end{longtable}

\begin{piege}
\texttt{==} compare les \emph{valeurs} ; \texttt{is} compare l'\emph{identite}
(meme objet en memoire). Utilise \texttt{is} uniquement avec \texttt{None} :
\texttt{if x is None}.
\end{piege}

% =============================================================================
\section{Chaines de caracteres}
% =============================================================================
Une chaine (\texttt{str}) est une suite ordonnee de caracteres. On y accede par
\emph{indexation} (\texttt{s[i]}, le premier caractere est a l'indice \texttt{0})
et par \emph{tranche} (\texttt{s[debut:fin]}, la borne de fin etant exclue). Les
indices negatifs comptent depuis la fin, ce qui evite de calculer la longueur.
Comme les chaines sont \textbf{immuables}, toutes les methodes qui « modifient »
une chaine (comme \texttt{upper} ou \texttt{replace}) renvoient en realite une
\emph{nouvelle} chaine : l'originale reste intacte. Les \emph{f-strings}
(\texttt{f"..."}) sont la facon moderne et lisible d'inserer des variables dans
un texte.

\begin{lstlisting}
s = "Python"
len(s)              # 6
s[0]                # 'P'      (indexation)
s[-1]               # 'n'      (indice negatif = depuis la fin)
s[0:3]              # 'Pyt'    (slicing [debut:fin])
s.upper()           # 'PYTHON'
s.replace("Py", "Ja")   # 'Jathon'
",".join(["a", "b"])    # 'a,b'
"a,b,c".split(",")      # ['a', 'b', 'c']

nom, age = "Alice", 30
f"{nom} a {age} ans"    # f-string : 'Alice a 30 ans'
\end{lstlisting}

\begin{piege}
Les chaines sont \emph{immuables} : \texttt{s[0] = "X"} leve
\texttt{TypeError}. Il faut reconstruire une nouvelle chaine.
\end{piege}

\subsection{Astuces de manipulation}
Au-dela des bases, une poignee de methodes reviennent sans cesse des qu'on doit
\emph{nettoyer}, \emph{aligner} ou \emph{analyser} du texte. Les connaitre evite
de reecrire a la main des boucles que la bibliotheque standard fait deja tres
bien.

\begin{lstlisting}
# Nettoyage : strip enleve les caracteres indesirables aux extremites
"  bonjour  ".strip()          # 'bonjour'   (espaces des deux cotes)
"##titre##".strip("#")         # 'titre'     (caractere precis)
"  bonjour".lstrip()           # 'bonjour'   (seulement a gauche)
"bonjour  ".rstrip()           # 'bonjour'   (seulement a droite)

# Alignement / padding : utile pour formater un affichage tabulaire
"7".zfill(3)                   # '007'  (complete avec des zeros a gauche)
"7".rjust(3, "0")              # '007'  (equivalent, plus general)
"ok".ljust(5, ".")             # 'ok...' (aligne a gauche)
"ok".center(6, "-")            # '-ok---' (centre)

# Recherche et tests
"algorithme".startswith("algo")   # True
"algorithme".endswith("me")       # True
"algorithme".find("go")           # 2   (-1 si absent, ne leve pas d'erreur)
"algorithme".index("go")          # 2   (leve ValueError si absent)
"abc123".isalpha()                # False (contient des chiffres)
"abc123".isalnum()                # True  (lettres + chiffres uniquement)

# Remplacement multiple d'un coup avec translate (rapide, sans boucle)
table = str.maketrans("ae", "AE")     # a->A, e->E
"algorithme".translate(table)         # 'AlgorithmE'

# split avec limite, et rsplit pour couper depuis la droite
"a=b=c".split("=", 1)             # ['a', 'b=c']       (une seule coupe)
"a/b/c.txt".rsplit("/", 1)        # ['a/b', 'c.txt']   (coupe depuis la fin)
"  a   b  c ".split()             # ['a', 'b', 'c']    (split() sans argument
                                   # ignore les espaces multiples/en trop)

# f-strings avancees : formater sans concatener
prix = 1234.5
f"{prix:,.2f}"          # '1,234.50'  (separateur de milliers, 2 decimales)
f"{prix:>10.1f}"        # '    1234.5' (aligne a droite sur 10 caracteres)
nom = "Alice"
f"{nom=}"                # "nom='Alice'" (affiche aussi le nom de variable, debug rapide)
\end{lstlisting}

\begin{infobox}{Compter des caracteres : le reflexe \texttt{Counter}}
Plutot que d'ecrire une boucle avec un dictionnaire pour compter les occurrences
de chaque caractere, \texttt{collections.Counter} le fait en une ligne et offre
directement les caracteres les plus frequents.
\begin{lstlisting}
from collections import Counter

freq = Counter("mississippi")      # Counter({'i': 4, 's': 4, 'p': 2, 'm': 1})
freq.most_common(2)                # [('i', 4), ('s', 4)]  (les 2 plus frequents)
freq["i"]                          # 4  (0 si absent, jamais de KeyError)
\end{lstlisting}
\end{infobox}

\begin{piege}
\texttt{"a" in "banane"} teste une \emph{sous-chaine}, pas un caractere isole
uniquement : \texttt{"an" in "banane"} vaut aussi \texttt{True}. Pratique, mais
ne pas confondre avec un test d'egalite caractere par caractere.
\end{piege}

% =============================================================================
\section{Structures de donnees}
% =============================================================================
Choisir la bonne structure de donnees est souvent la moitie de la solution d'un
probleme. Python en offre quatre fondamentales, qu'on distingue selon trois
questions : les elements sont-ils \emph{ordonnes} ? la collection est-elle
\emph{modifiable} ? les acces se font-ils par position ou par cle ? Un bon
reflexe : si tu ecris une boucle pour verifier « est-ce que x est dans la
liste ? » de facon repetee, un \texttt{set} ou un \texttt{dict} (recherche en
\texttt{O(1)} moyen) remplace avantageusement la liste (recherche en
\texttt{O(n)}).

\subsection{Liste -- ordonnee, modifiable}
La liste est la structure a tout faire : une sequence ordonnee que l'on peut
agrandir, trier et modifier sur place. Les \emph{comprehensions} de liste sont
une facon concise et rapide de construire une liste a partir d'une autre
sequence, en remplacant une boucle \texttt{for} et des \texttt{append}.
\begin{lstlisting}
nums = [3, 1, 2]
nums.append(4)      # [3, 1, 2, 4]
nums.sort()         # [1, 2, 3, 4]  (sur place)
nums[1:3]           # [2, 3]        (slicing)
carres = [x*x for x in range(5)]    # comprehension : [0,1,4,9,16]
\end{lstlisting}

\subsection{Tuple -- ordonne, immuable}
Le tuple ressemble a une liste mais ne peut plus changer une fois cree. On
l'utilise pour des donnees qui vont ensemble et ne doivent pas bouger (un point
\texttt{(x, y)}, une ligne d'une base de donnees) et pour le \emph{deballage}
(unpacking) qui affecte plusieurs variables en une ligne. Etant immuable, un
tuple peut servir de cle de dictionnaire, contrairement a une liste.

\begin{lstlisting}
point = (3, 4)
x, y = point        # deballage (unpacking)
# point[0] = 9      -> TypeError : un tuple ne se modifie pas
\end{lstlisting}

\subsection{Dictionnaire -- cle -> valeur}
Le dictionnaire associe des \emph{cles} a des \emph{valeurs} et permet de
retrouver une valeur instantanement a partir de sa cle. C'est l'outil ideal pour
compter, indexer, ou representer un objet par ses attributs. La methode
\texttt{get(cle, defaut)} evite un \texttt{KeyError} lorsque la cle peut etre
absente. Depuis Python 3.7, un dictionnaire conserve l'ordre d'insertion des
cles.

\begin{lstlisting}
age = {"Alice": 30, "Bob": 25}
age["Alice"]            # 30
age.get("Zoe", 0)       # 0  (valeur par defaut si absente)
age["Zoe"] = 22         # ajout
for cle, val in age.items():
    print(cle, val)
\end{lstlisting}

\subsubsection*{Astuces de manipulation}
Le dictionnaire est la structure la plus polyvalente de Python : elle sert a la
fois de \emph{table de correspondance}, de \emph{compteur} et d'\emph{index}.
Quelques idiomes reviennent dans presque tous les problemes d'algorithmique.

\begin{lstlisting}
# Dict comprehension : construire un dict en une ligne, comme pour une liste
carres = {x: x*x for x in range(5)}          # {0:0, 1:1, 2:4, 3:9, 4:16}
mots = ["chat", "chien", "oiseau"]
longueurs = {m: len(m) for m in mots}        # {'chat':4, 'chien':5, 'oiseau':7}

# setdefault : recupere une cle, ou l'insere avec une valeur par defaut
graphe = {}
graphe.setdefault("A", []).append("B")       # cree [] si absent, puis append
graphe.setdefault("A", []).append("C")       # {'A': ['B', 'C']}

# defaultdict : evite setdefault repete, la valeur par defaut est automatique
from collections import defaultdict
groupes = defaultdict(list)
for mot in ["chat", "chien", "chameau"]:
    groupes[mot[0]].append(mot)              # pas besoin de tester si 'c' existe
# {'c': ['chat', 'chien', 'chameau']}

# Fusionner deux dicts (Python 3.9+) : | cree un nouveau dict, |= modifie en place
defauts = {"theme": "clair", "langue": "fr"}
utilisateur = {"langue": "en"}
config = defauts | utilisateur               # {'theme':'clair', 'langue':'en'}
                                              # (les cles de droite ecrasent celles de gauche)

# Parcourir cles / valeurs / paires
list(age.keys())        # ['Alice', 'Bob', 'Zoe']
list(age.values())      # [30, 25, 22]
list(age.items())       # [('Alice', 30), ('Bob', 25), ('Zoe', 22)]

# Trier un dict par valeur (frequent : classement par score, par frequence...)
scores = {"Alice": 90, "Bob": 75, "Zoe": 88}
sorted(scores.items(), key=lambda kv: kv[1], reverse=True)
# [('Alice', 90), ('Zoe', 88), ('Bob', 75)]  (trie decroissant sur la valeur)

# Inverser un dict (valeur -> cle), valable seulement si les valeurs sont uniques
{v: k for k, v in age.items()}

# Supprimer une cle sans planter si elle est absente
age.pop("Inconnu", None)     # None si absente ; pas de KeyError
\end{lstlisting}

\begin{piege}
Une cle de dictionnaire doit etre \emph{immuable} (\texttt{hashable}) : \texttt{str},
\texttt{int}, \texttt{tuple} conviennent, mais une \texttt{list} comme cle leve
\texttt{TypeError: unhashable type}. Utilise un \texttt{tuple} si tu as besoin
d'une cle composee, par exemple \texttt{grille[(x, y)] = valeur}.
\end{piege}

\subsection{Ensemble (set) -- valeurs uniques}
Un ensemble ne contient jamais de doublon et n'a pas d'ordre. On l'emploie pour
eliminer les repetitions (\texttt{set(ma\_liste)}), tester tres vite une
appartenance, ou faire des operations mathematiques : union (\texttt{|}),
intersection (\texttt{\&}), difference (\texttt{-}). C'est souvent l'astuce qui
transforme un algorithme lent en algorithme rapide.

\begin{lstlisting}
vus = {1, 2, 2, 3}      # {1, 2, 3}
vus.add(4)
3 in vus                # True  (appartenance en O(1) moyen)
{1, 2} & {2, 3}         # {2}   (intersection)
\end{lstlisting}

\begin{infobox}{Quel type choisir ?}
Ordre + modifications $\rightarrow$ \textbf{liste}. Donnees fixes $\rightarrow$
\textbf{tuple}. Correspondance cle/valeur $\rightarrow$ \textbf{dict}. Unicite et
appartenance rapide $\rightarrow$ \textbf{set}.
\end{infobox}

% =============================================================================
\section{Conditions et boucles}
% =============================================================================
Les conditions (\texttt{if}/\texttt{elif}/\texttt{else}) et les boucles
(\texttt{for}, \texttt{while}) controlent le \emph{flux} d'execution : quelles
instructions s'executent, et combien de fois. En Python, les blocs ne sont pas
delimites par des accolades mais par l'\textbf{indentation} : tout ce qui est
decale d'un cran vers la droite appartient au bloc. Une indentation incoherente
provoque une \texttt{IndentationError}. La boucle \texttt{for} parcourt
directement les elements d'une sequence (pas des indices) ; \texttt{enumerate}
fournit en prime la position, et \texttt{range} genere une suite d'entiers.

\begin{lstlisting}
# Condition
if n > 0:
    signe = "positif"
elif n == 0:
    signe = "nul"
else:
    signe = "negatif"

# Boucle for
for i in range(3):          # 0, 1, 2
    print(i)

for i, val in enumerate(["a", "b"]):   # (0,'a') (1,'b')
    print(i, val)

# Boucle while
while n > 0:
    n -= 1

# break / continue / else
for x in nums:
    if x == cible:
        break               # sort de la boucle
else:
    print("non trouve")     # execute si pas de break
\end{lstlisting}

\begin{piege}
\texttt{range(a, b)} s'arrete a \texttt{b-1} (borne haute exclue). \texttt{range(5)}
produit \texttt{0,1,2,3,4}.
\end{piege}

% =============================================================================
\section{Fonctions}
% =============================================================================
Une fonction regroupe un morceau de code reutilisable sous un nom. Elle prend des
\emph{parametres} en entree et renvoie un resultat avec \texttt{return} (sans
\texttt{return}, elle renvoie \texttt{None}). Bien decouper un programme en
petites fonctions au role unique le rend plus lisible, plus facile a tester et a
deboguer. Python offre des arguments \emph{par defaut}, des arguments
\emph{nommes} (qui clarifient les appels), et la syntaxe \texttt{*args} /
\texttt{**kwargs} pour accepter un nombre variable d'arguments.

\begin{lstlisting}
def salue(nom, salutation="Bonjour"):   # argument par defaut
    return f"{salutation}, {nom} !"

salue("Alice")                 # 'Bonjour, Alice !'
salue("Bob", "Salut")          # 'Salut, Bob !'
salue(nom="Zoe")               # argument nomme

# Nombre variable d'arguments
def total(*args, **kwargs):    # args = tuple, kwargs = dict
    return sum(args)

total(1, 2, 3)                 # 6

# Fonction anonyme (lambda)
carre = lambda x: x * x
sorted(mots, key=lambda m: len(m))     # trier par longueur
\end{lstlisting}

\begin{piege}
Ne jamais utiliser un objet \emph{mutable} comme valeur par defaut :
\texttt{def f(x, liste=[])} partage la meme liste entre les appels. Ecris
\texttt{def f(x, liste=None): liste = liste or []}.
\end{piege}

% =============================================================================
\section{Portee des variables (scope)}
% =============================================================================
\begin{lstlisting}
compteur = 0                   # variable globale

def incr():
    global compteur            # sinon : variable locale
    compteur += 1

def exterieure():
    x = 10
    def interieure():
        nonlocal x             # modifie x de la fonction englobante
        x += 1
    interieure()
    return x                   # 11
\end{lstlisting}

\begin{infobox}{Regle LEGB}
Python cherche un nom dans cet ordre : \textbf{L}ocal $\rightarrow$
\textbf{E}nglobant $\rightarrow$ \textbf{G}lobal $\rightarrow$ \textbf{B}uilt-in.
\end{infobox}

% =============================================================================
\section{Classes et objets}
% =============================================================================
Une \emph{classe} est un modele (un plan) qui decrit des donnees (les
\emph{attributs}) et des comportements (les \emph{methodes}). Un \emph{objet} est
une instance concrete de ce modele. Le premier parametre de chaque methode,
\texttt{self}, designe l'objet sur lequel on travaille ; la methode speciale
\texttt{\_\_init\_\_} est le \emph{constructeur}, appele automatiquement a la
creation. La \emph{programmation orientee objet} permet de regrouper l'etat et le
comportement au meme endroit, et l'\emph{heritage} (via \texttt{super()}) de
reutiliser et specialiser une classe existante sans la reecrire.

\begin{lstlisting}
class Compte:
    taux = 0.02                     # attribut de classe (partage)

    def __init__(self, solde):      # constructeur
        self.solde = solde          # attribut d'instance

    def deposer(self, montant):     # methode
        self.solde += montant

    def __str__(self):              # representation lisible
        return f"Compte({self.solde} EUR)"

c = Compte(100)
c.deposer(50)
print(c)                            # Compte(150 EUR)

# Heritage
class ComptePremium(Compte):
    def __init__(self, solde, bonus):
        super().__init__(solde)     # appelle le parent
        self.bonus = bonus
\end{lstlisting}

\begin{piege}
Oublier \texttt{self} en premier parametre des methodes leve
\texttt{TypeError: ... takes 0 positional arguments but 1 was given}.
\end{piege}

% =============================================================================
\section{Gestion des erreurs (exceptions)}
% =============================================================================
Une erreur non geree arrete le programme et affiche une \emph{traceback}. Plutot
que de la subir, on peut l'\emph{anticiper} : le bloc \texttt{try} contient le
code risque, et chaque \texttt{except} rattrape un type d'erreur precis pour
reagir proprement. C'est la difference entre un programme qui plante brutalement
et un programme robuste qui affiche un message clair et continue. Le bloc
\texttt{else} s'execute si aucune erreur n'est survenue, et \texttt{finally}
s'execute \emph{toujours} (succes ou echec) -- ideal pour liberer une ressource,
fermer un fichier ou une connexion. On peut aussi \emph{declencher} soi-meme une
erreur avec \texttt{raise} quand une donnee est invalide.

\begin{lstlisting}
try:
    x = int(input("Nombre : "))
    resultat = 10 / x
except ValueError:                 # conversion impossible
    print("Ce n'est pas un entier.")
except ZeroDivisionError:          # division par zero
    print("Division par zero interdite.")
except Exception as e:             # filet de securite
    print(f"Erreur inattendue : {e}")
else:
    print("Tout s'est bien passe :", resultat)
finally:
    print("Toujours execute (nettoyage).")

# Lever sa propre exception
def racine(x):
    if x < 0:
        raise ValueError("x doit etre positif")
    return x ** 0.5
\end{lstlisting}

\begin{piege}
Evite \texttt{except:} tout seul (attrape meme \texttt{Ctrl+C}). Attrape
toujours le type le plus precis possible, et place les \texttt{except} du plus
specifique au plus general.
\end{piege}

% =============================================================================
\section{Les erreurs Python les plus frequentes}
% =============================================================================
Savoir \emph{lire} une erreur fait gagner un temps precieux : le nom de
l'exception dit \emph{quel genre} de probleme s'est produit, et le message dit
\emph{ou} et \emph{pourquoi}. Voici les exceptions que tu rencontreras le plus
souvent en debutant. Prends l'habitude d'associer chaque nom a sa cause typique :
avec un peu de pratique, tu diagnostiqueras la plupart des bugs en lisant
simplement la premiere ligne du message.

\renewcommand{\arraystretch}{1.35}
\begin{longtable}{>{\raggedright\arraybackslash}p{4.2cm}
                  >{\raggedright\arraybackslash}p{8.6cm}}
\rowcolor{primary}
\textbf{\color{white}Exception} & \textbf{\color{white}Cause typique} \\
\hline
\texttt{SyntaxError}       & Faute de syntaxe (\texttt{:} oublie, parenthese non fermee). \\
\texttt{IndentationError}  & Indentation incorrecte ou incoherente (espaces vs tabulations). \\
\texttt{NameError}         & Variable utilisee avant d'etre definie. \\
\texttt{TypeError}         & Operation entre types incompatibles (\texttt{"a" + 1}). \\
\texttt{ValueError}        & Bon type mais valeur invalide (\texttt{int("abc")}). \\
\texttt{IndexError}        & Indice hors des bornes d'une liste. \\
\texttt{KeyError}          & Cle absente d'un dictionnaire. \\
\texttt{AttributeError}    & Attribut/methode inexistant sur un objet. \\
\texttt{ZeroDivisionError} & Division ou modulo par zero. \\
\texttt{ImportError}       & Module ou nom introuvable a l'import. \\
\texttt{RecursionError}    & Recursion trop profonde (cas de base manquant). \\
\end{longtable}

\begin{infobox}{Lire une traceback}
Lis-la \textbf{de bas en haut} : la derniere ligne donne le type d'erreur et le
message ; remonte pour trouver le fichier et la ligne exacte ou elle s'est
produite.
\end{infobox}

% =============================================================================
\section{Modules et imports}
% =============================================================================
Un \emph{module} est un fichier Python que l'on peut reutiliser dans un autre
grace a \texttt{import}. La \emph{bibliotheque standard} fournit deja des
centaines de modules prets a l'emploi (\texttt{math}, \texttt{collections},
\texttt{datetime}, \texttt{random}\ldots) -- inutile de reinventer la roue. On
peut importer un module entier, seulement certains noms (\texttt{from ... import
...}), ou lui donner un alias plus court (\texttt{import numpy as np}). Le test
\texttt{if \_\_name\_\_ == "\_\_main\_\_":} permet qu'un fichier serve a la fois
de programme executable et de module importable.

\begin{lstlisting}
import math
math.sqrt(16)                      # 4.0

from collections import Counter, defaultdict
Counter("aab")                     # Counter({'a': 2, 'b': 1})

import datetime as dt              # alias
dt.date.today()

# Point d'entree d'un script
if __name__ == "__main__":
    main()                         # execute seulement si lance directement
\end{lstlisting}

% =============================================================================
\section{Iterateurs et generateurs}
% =============================================================================
Un \emph{iterateur} est un objet que l'on peut parcourir element par element avec
une boucle \texttt{for}. Un \emph{generateur} est la facon la plus simple d'en
creer un : c'est une fonction qui, au lieu de \texttt{return}, utilise
\texttt{yield} pour livrer ses valeurs \emph{une a une}, en se mettant en pause
entre chaque. Cette approche \emph{paresseuse} (lazy) ne calcule la valeur
suivante qu'au moment ou on en a besoin : la memoire consommee reste constante,
meme pour parcourir des millions d'elements ou un flux potentiellement infini.

\begin{lstlisting}
def compte_jusqua(n):
    i = 0
    while i < n:
        yield i                    # renvoie une valeur et se met en pause
        i += 1

for v in compte_jusqua(3):         # 0, 1, 2
    print(v)

# Expression generatrice (parentheses au lieu de crochets)
somme = sum(x*x for x in range(1000))   # ne cree pas la liste intermediaire
\end{lstlisting}

\begin{infobox}{Interet}
Ideal pour de gros flux de donnees : la memoire reste constante quel que soit le
nombre d'elements parcourus.
\end{infobox}

% =============================================================================
\section{Decorateurs}
% =============================================================================
Un decorateur est une fonction qui \emph{enveloppe} une autre fonction pour lui
ajouter un comportement (journalisation, mesure de temps, mise en cache,
verification des droits\ldots) \textbf{sans toucher a son code}. C'est possible
parce qu'en Python les fonctions sont des \emph{objets de premiere classe} : on
peut les passer en argument, les renvoyer, et les stocker dans des variables. La
syntaxe \texttt{@decorateur} placee au-dessus d'une fonction est simplement un
raccourci pour \texttt{ma\_fonction = decorateur(ma\_fonction)}. Le decorateur
renvoie une fonction \emph{wrapper} qui execute du code avant et/ou apres la
fonction d'origine. On l'utilise partout ou l'on veut appliquer le meme
traitement a plusieurs fonctions sans dupliquer ce code.

\begin{lstlisting}
import functools
import time

def chronometre(func):
    @functools.wraps(func)             # conserve nom / docstring
    def wrapper(*args, **kwargs):
        debut = time.perf_counter()
        resultat = func(*args, **kwargs)
        duree = time.perf_counter() - debut
        print(f"{func.__name__} a pris {duree:.4f}s")
        return resultat
    return wrapper

@chronometre                            # equivaut a : lent = chronometre(lent)
def lent():
    time.sleep(0.1)

lent()      # affiche : lent a pris 0.10xx s
\end{lstlisting}

\subsection{Decorateurs integres utiles}
\begin{lstlisting}
from functools import lru_cache

@lru_cache(maxsize=None)                # memoise les resultats
def fib(n):
    return n if n < 2 else fib(n-1) + fib(n-2)

class Cercle:
    def __init__(self, r):
        self._r = r

    @property                           # acces comme un attribut
    def aire(self):
        return 3.14159 * self._r ** 2

    @staticmethod                       # pas de self
    def unite():
        return "cm"
\end{lstlisting}

\begin{piege}
Sans \texttt{@functools.wraps}, la fonction decoree perd son nom et sa docstring
(\texttt{func.\_\_name\_\_} devient \texttt{'wrapper'}), ce qui gene le debogage.
\end{piege}

\newpage
% =============================================================================
\section{Fiche recapitulative}
% =============================================================================
\renewcommand{\arraystretch}{1.35}
\begin{longtable}{>{\raggedright\arraybackslash}p{4.5cm}
                  >{\raggedright\arraybackslash}p{8.3cm}}
\rowcolor{primary}
\textbf{\color{white}Concept} & \textbf{\color{white}A retenir} \\
\hline
Typage            & Dynamique et fort ; \texttt{type(x)}, \texttt{isinstance(x, T)}. \\
\texttt{/} vs \texttt{//} & \texttt{/} = float ; \texttt{//} = division entiere. \\
Immuable          & \texttt{str}, \texttt{tuple}, \texttt{int} ne se modifient pas. \\
Comprehension     & \texttt{[f(x) for x in it if cond]}. \\
\texttt{==} vs \texttt{is} & valeur vs identite ; \texttt{is} avec \texttt{None}. \\
Defaut mutable    & Ne jamais mettre \texttt{[]} ou \texttt{\{\}} en defaut. \\
Exceptions        & \texttt{try/except/else/finally} ; du specifique au general. \\
Traceback         & Se lit de bas en haut. \\
Generateur        & \texttt{yield} : valeurs paresseuses, memoire constante. \\
Decorateur        & Enveloppe une fonction ; \texttt{@functools.wraps}. \\
\texttt{lru\_cache}   & Memoisation automatique des appels. \\
\texttt{\_\_main\_\_}   & \texttt{if \_\_name\_\_ == "\_\_main\_\_":} pour le point d'entree. \\
\end{longtable}

\vspace{0.6cm}
\begin{infobox}{Methode pour deboguer un probleme}
\begin{enumerate}[leftmargin=*,itemsep=2pt,topsep=2pt]
  \item Lire le \textbf{type} d'exception et le message (derniere ligne).
  \item Localiser le \textbf{fichier} et la \textbf{ligne} dans la traceback.
  \item Afficher les valeurs suspectes (\texttt{print} ou \texttt{breakpoint()}).
  \item Reproduire avec un \textbf{petit} exemple minimal.
  \item Corriger, puis verifier les \textbf{cas limites} (vide, zero, negatif).
\end{enumerate}
\end{infobox}

\end{document}
